\subsection{Affine Maps and Affine Groups}

Throughout this subsection, $\mathcal{V}$ is finite-dimensional over a fixed field $\mathbb{F}$.

\begin{definition} \label{def:affine-transformation} \hlt{Affine Transformation}\\
A map $f:\mathcal{V}\to\mathcal{V}$ is \hlt{affine} iff $(\forall\,n\in\mathbb{Z}_{>0})(\forall\,(\mathbf{x}_1,\ldots,\mathbf{x}_n)\in\mathcal{V}^n)(\forall\,\mathbf{r}\in\mathbb{F}^n)\,[\sum\limits_{j=1}^n r_j=1\Rightarrow f(\sum\limits_{j=1}^n r_j\mathbf{x}_j)=\sum\limits_{j=1}^n r_jf(\mathbf{x}_j)]$.
\end{definition}

\begin{theorem} \label{thm:affine-map-two-point-test} \hlt{Two-Point Test for Affine Maps}\\
Assume $\operatorname{char}(\mathbb{F})\neq2$. A map $f:\mathcal{V}\to\mathcal{V}$ is affine iff $(\forall\,\mathbf{x},\mathbf{y}\in\mathcal{V})(\forall\,r\in\mathbb{F})\,f(r\mathbf{x}+(1-r)\mathbf{y})=rf(\mathbf{x})+(1-r)f(\mathbf{y})$.
\begin{proof}
\begin{enumerate}[label=\roman*.]
\setlength{\itemsep}{0pt}
\item[$(\Rightarrow)$] The identity holds by Definition \ref{def:affine-transformation} with $n=2$.
\item[$(\Leftarrow)$] Let $\Gamma_f\defeq\{(\mathbf{x},f(\mathbf{x})):\mathbf{x}\in\mathcal{V}\}\subseteq\mathcal{V}\times\mathcal{V}$. $\Gamma_f$ is two-affine closed as $r(\mathbf{x},f(\mathbf{x}))+(1-r)(\mathbf{y},f(\mathbf{y}))=(r\mathbf{x}+(1-r)\mathbf{y},f(r\mathbf{x}+(1-r)\mathbf{y}))$.\\
The nonempty set $\Gamma_f$ is affine closed by Theorem \ref{thm:affine-closed-two-affine-closed} in $\mathcal{V}\times\mathcal{V}$. Fix $n\in\mathbb{Z}_{>0}$, $(\mathbf{x}_1,\ldots,\mathbf{x}_n)\in\mathcal{V}^n$, and $\mathbf{r}\in\mathbb{F}^n$ with $\sum\limits_{j=1}^n r_j=1$.\\
Then $(\sum\limits_{j=1}^n r_j\mathbf{x}_j,\sum\limits_{j=1}^n r_jf(\mathbf{x}_j))\in\Gamma_f$, so $f(\sum\limits_{j=1}^n r_j\mathbf{x}_j)=\sum\limits_{j=1}^n r_jf(\mathbf{x}_j)$. Thus $f$ is affine.
\end{enumerate}
\hlt{Structural Commentary}: Convert binary preservation into two-affine closure of the graph, then invoke affine closure.\\
\hlt{Key Insight}: Affine maps are exactly the maps with affine-closed graphs.
\end{proof}
\end{theorem}

\begin{definition} \label{def:translation-map} \hlt{Translations}\\
For $\mathbf{a}\in\mathcal{V}$, the \hlt{translation by $\mathbf{a}$} is the map $T_{\mathbf{a}}:\mathcal{V}\to\mathcal{V}$ defined by $T_{\mathbf{a}}(\mathbf{x})=\mathbf{x}+\mathbf{a}$.
\end{definition}

\begin{theorem} \label{thm:affine-map-linear-part-translation} \hlt{Affine Maps Are Linear Maps Followed by Translations}\\
A map $f:\mathcal{V}\to\mathcal{V}$ is affine iff $(\exists!\,(\mathbf{a},\tau)\in\mathcal{V}\times\LL(\mathcal{V}))\,f=T_{\mathbf{a}}\circ\tau$. The unique pair satisfies $\mathbf{a}=f(\mathbf{0})$ and $\tau=T_{-\mathbf{a}}\circ f$.
\begin{proof}
\begin{enumerate}[label=\roman*.] 
\setlength{\itemsep}{0pt}
\item[$(\Rightarrow)$] Put $\mathbf{a}\defeq f(\mathbf{0})$ and $\tau\defeq T_{-\mathbf{a}}\circ f$. For $r,s\in\mathbb{F}$ and $\mathbf{x},\mathbf{y}\in\mathcal{V}$, the identity $\tau(r\mathbf{x}+s\mathbf{y})=rf(\mathbf{x})+sf(\mathbf{y})+(1-r-s)\mathbf{a}-\mathbf{a}=r\tau(\mathbf{x})+s\tau(\mathbf{y})$ holds by affine preservation with coefficients $r,s,1-r-s$. Hence $\tau\in\LL(\mathcal{V})$ and $f=T_{\mathbf{a}}\circ\tau$.\\
If $f=T_{\mathbf{b}}\circ\sigma$ with $(\mathbf{b},\sigma)\in\mathcal{V}\times\LL(\mathcal{V})$, then $\mathbf{b}=f(\mathbf{0})=\mathbf{a}$ and $\sigma=T_{-\mathbf{b}}\circ f=\tau$, proving uniqueness.
\item[$(\Leftarrow)$] Suppose $f=T_{\mathbf{a}}\circ\tau$ with $\tau\in\LL(\mathcal{V})$. Fix $n\in\mathbb{Z}_{>0}$, $(\mathbf{x}_1,\ldots,\mathbf{x}_n)\in\mathcal{V}^n$, and $\mathbf{r}\in\mathbb{F}^n$ with $\sum\limits_{j=1}^n r_j=1$.\\
Then $f(\sum\limits_{j=1}^n r_j\mathbf{x}_j)=\mathbf{a}+\sum\limits_{j=1}^n r_j\tau(\mathbf{x}_j)=\sum\limits_{j=1}^n r_j(\mathbf{a}+\tau(\mathbf{x}_j))=\sum\limits_{j=1}^n r_jf(\mathbf{x}_j)$, so $f$ is affine.
\end{enumerate}
\hlt{Structural Commentary}: Translate $f(\mathbf{0})$ to zero, verify linearity, then recover and identify the factorization.\\
\hlt{Key Insight}: The image of zero determines the entire translation part.
\end{proof}
\end{theorem}

\begin{definition} \label{def:affine-group-translations-linear-part} \hlt{Affine Group Data}
\begin{enumerate}[label=(\roman*)]
\setlength{\itemsep}{0pt}
\item The \hlt{general linear group} is $\operatorname{GL}(\mathcal{V})\defeq\{\tau\in\LL(\mathcal{V}):(\exists\,\sigma\in\LL(\mathcal{V}))\,[\sigma\tau=I_{\mathcal{V}}=\tau\sigma]\}$.
\item The \hlt{affine group} is $\operatorname{Aff}(\mathcal{V})\defeq\{T_{\mathbf{a}}\circ\tau:\mathbf{a}\in\mathcal{V},\,\tau\in\operatorname{GL}(\mathcal{V})\}$.
\item The \hlt{translation group} is $\operatorname{Trans}(\mathcal{V})\defeq\{T_{\mathbf{a}}:\mathbf{a}\in\mathcal{V}\}$.
\item The \hlt{linear-part map} is $\Lambda:\operatorname{Aff}(\mathcal{V})\to\operatorname{GL}(\mathcal{V})$, defined by $\Lambda(T_{\mathbf{a}}\circ\tau)=\tau$; it is well defined by uniqueness in Theorem \ref{thm:affine-map-linear-part-translation}.
\end{enumerate}
\end{definition}

\begin{corollary} \label{cor:affine-map-bijective-linear-part} \hlt{Bijective Affine Maps}\\
Let $\mathbf{a}\in\mathcal{V}$ and $\tau\in\LL(\mathcal{V})$. Then $T_{\mathbf{a}}\circ\tau$ is bijective iff $\tau\in\operatorname{GL}(\mathcal{V})$.
\begin{proof}
\begin{enumerate}[label=\roman*.]
\setlength{\itemsep}{0pt}
\item[$(\Rightarrow)$] The translation $T_{-\mathbf{a}}$ is bijective with inverse $T_{\mathbf{a}}$. Hence $\tau=T_{-\mathbf{a}}\circ(T_{\mathbf{a}}\circ\tau)$ is bijective, so $\tau\in\operatorname{GL}(\mathcal{V})$ by Proposition \ref{prop:linear-map-invertible-iff-bijective}.
\item[$(\Leftarrow)$] If $\tau\in\operatorname{GL}(\mathcal{V})$, then $T_{\mathbf{a}}\circ\tau$ is bijective with inverse $\tau^{-1}\circ T_{-\mathbf{a}}$.
\end{enumerate}
\hlt{Structural Commentary}: Cancel the translation in one direction and compose the two explicit inverses in the other.\\
\hlt{Key Insight}: Translation does not affect bijectivity.
\end{proof}
\end{corollary}

\begin{corollary} \label{cor:affine-group-linear-part} \hlt{Affine Group and Translations}
\begin{enumerate}[label=(\roman*)]
\setlength{\itemsep}{0pt}
\item $(\forall\,\mathbf{a},\mathbf{b}\in\mathcal{V})(\forall\,\tau,\sigma\in\LL(\mathcal{V}))\,(T_{\mathbf{a}}\circ\tau)\circ(T_{\mathbf{b}}\circ\sigma)=T_{\mathbf{a}+\tau(\mathbf{b})}\circ(\tau\sigma)$.
\item $\operatorname{GL}(\mathcal{V})$ is a group under composition.
\item $\operatorname{Aff}(\mathcal{V})$ is a group under composition, with identity $T_{\mathbf{0}}\circ I_{\mathcal{V}}$ and $(T_{\mathbf{a}}\circ\tau)^{-1}=T_{-\tau^{-1}(\mathbf{a})}\circ\tau^{-1}$.
\item $\operatorname{Trans}(\mathcal{V})\trianglelefteq\operatorname{Aff}(\mathcal{V})$.
\item $\Lambda$ is a surjective group homomorphism and $\ker\Lambda=\operatorname{Trans}(\mathcal{V})$.
\item $\operatorname{Aff}(\mathcal{V})/\operatorname{Trans}(\mathcal{V})\cong\operatorname{GL}(\mathcal{V})$.
\end{enumerate}
\begin{proof}
\begin{enumerate}[label=(\roman*)]
\setlength{\itemsep}{0pt}
\item For $\mathbf{x}\in\mathcal{V}$, $((T_{\mathbf{a}}\circ\tau)\circ(T_{\mathbf{b}}\circ\sigma))(\mathbf{x})=\mathbf{a}+\tau(\mathbf{b})+(\tau\sigma)(\mathbf{x})$, which proves the map identity.
\item Composition is associative and $I_{\mathcal{V}}\in\operatorname{GL}(\mathcal{V})$. Closure under products holds by Proposition \ref{prop:inverse-of-product-linear-maps}(i).\\
If $\tau\in\operatorname{GL}(\mathcal{V})$, then $\tau^{-1}\in\LL(\mathcal{V})$ with inverse $\tau$ by Definition \ref{def:invertible-linear-map-and-isomorphism} and Proposition \ref{prop:inverse-linear-map-is-unique}, so $\tau^{-1}\in\operatorname{GL}(\mathcal{V})$.
\item Closure holds by (i)--(ii), and composition is associative. The identity is $T_{\mathbf{0}}\circ I_{\mathcal{V}}$.\\
By (i), $T_{-\tau^{-1}(\mathbf{a})}\circ\tau^{-1}$ composes with $T_{\mathbf{a}}\circ\tau$ in both orders to $T_{\mathbf{0}}\circ I_{\mathcal{V}}$, proving the inverse formula and the group claim.
\item The identities $T_{\mathbf{a}}T_{\mathbf{b}}=T_{\mathbf{a}+\mathbf{b}}$, $T_{\mathbf{0}}=I_{\mathcal{V}}$, and $T_{\mathbf{a}}^{-1}=T_{-\mathbf{a}}$ hold by (i), so $\operatorname{Trans}(\mathcal{V})$ is a subgroup.\\
If $g=T_{\mathbf{a}}\circ\tau$, then $gT_{\mathbf{b}}g^{-1}=T_{\tau(\mathbf{b})}$ by (i). Since $\tau$ is bijective by Proposition \ref{prop:linear-map-invertible-iff-bijective}, $g\operatorname{Trans}(\mathcal{V})g^{-1}=\operatorname{Trans}(\mathcal{V})$.
\item By (i), $\Lambda((T_{\mathbf{a}}\circ\tau)(T_{\mathbf{b}}\circ\sigma))=\tau\sigma=\Lambda(T_{\mathbf{a}}\circ\tau)\Lambda(T_{\mathbf{b}}\circ\sigma)$, so $\Lambda$ is a homomorphism.\\
It is surjective because $\Lambda(T_{\mathbf{0}}\circ\tau)=\tau$, and $\ker\Lambda=\{T_{\mathbf{a}}\circ I_{\mathcal{V}}:\mathbf{a}\in\mathcal{V}\}=\operatorname{Trans}(\mathcal{V})$.
\item Put $N\defeq\operatorname{Trans}(\mathcal{V})$. The rule $\overline{\Lambda}(gN)=\Lambda(g)$ defines $\overline{\Lambda}:\operatorname{Aff}(\mathcal{V})/N\to\operatorname{GL}(\mathcal{V})$ as $gN=hN\Rightarrow h^{-1}g\in N=\ker\Lambda\Rightarrow\Lambda(g)=\Lambda(h)$.\\
It is a surjective homomorphism by (v), and $\ker\overline{\Lambda}=\{N\}$ because $\overline{\Lambda}(gN)=I_{\mathcal{V}}\iff g\in\ker\Lambda=N\iff gN=N$. QED.
\end{enumerate}
\hlt{Structural Commentary}: Derive product and inverse laws, prove normality, compute the kernel, then descend $\Lambda$ to the quotient.\\
\hlt{Key Insight}: The linear-part map forgets exactly the translation subgroup.
\end{proof}
\end{corollary}
